% Generated draft: requires mathematical and source review.
\documentclass[11pt]{article}
\usepackage[margin=25mm]{geometry}
\usepackage{fontspec}
\setmainfont{Latin Modern Roman}
\usepackage{amsmath,amssymb,mathtools,unicode-math}
\setmathfont{Latin Modern Math}
\usepackage{enumitem,needspace,xurl,hyperref}
\setlength{\parindent}{0pt}
\setlength{\parskip}{5pt}
\setlength{\emergencystretch}{3em}
\newcommand{\N}{\mathbb{N}}
\newcommand{\Z}{\mathbb{Z}}
\newcommand{\Q}{\mathbb{Q}}
\newcommand{\R}{\mathbb{R}}
\newcommand{\C}{\mathbb{C}}
\newcommand{\F}{\mathbb{F}}
\newcommand{\A}{\mathbb{A}}
\newcommand{\PP}{\mathbb{P}}
\newcommand{\E}{\mathbb{E}}
\newcommand{\eps}{\varepsilon}
\newcommand{\dd}{\,\mathrm{d}}
\begin{document}
\section*{UnsolvedMath ID: 20000317}
\textbf{Source problem code:} catalogue entry
\subsection{Definitions and mathematical statement}
Cheltsov asks, for a smooth Fano variety $X$ with $\operatorname{Pic}X\simeq\Z$, whether birational superrigidity is equivalent to
\[
\forall m\in\Z_{>0}\ \forall\text{ movable }\mathcal M\subset|-mK_X|:
\quad(X,\tfrac1m\mathcal M)\text{ is canonical}.\tag{1}
\]
Movable means having no fixed divisorial components. Birational superrigidity is understood in the Mori-fiber-space category, requiring birational maps from the rank-one Fano variety to other Mori fiber spaces to be biregular in the appropriate sense. The original attributes the superrigidity-to-(1) implication to MMP.

It also introduces the two distinct readings
\[
\forall m>0\ \forall D\in|-mK_X|:\quad
(X,\tfrac1mD)\text{ is canonical},\tag{2a}
\]
or the same condition with ``purely log terminal'' in place of ``canonical,'' denoted (2b). These are singularity conditions on all individual divisors, stronger in scope than testing movable systems; the canonical and plt alternatives must not be conflated. Canonicity is tested by nonnegative exceptional discrepancies, with the usual boundary conventions. The source reports that sufficiently general degree-$N$ hypersurfaces $X_N\subset\PP^N$ satisfy (1) and (2) for $N\ge6$, attributing this to Pukhlikov [8], and asks what happens for $N=4,5$.

The attached unreviewed aid reports that the Noether--Fano criterion answers the equivalence and that de Fernex's theorem with its erratum yields (1) for every smooth quartic threefold and quintic fourfold. It leaves the general quartic/quintic (2a) and (2b) questions unresolved in its literature assessment. For each $N\ge4$ it proposes a special smooth degree-$N$ hypersurface satisfying (1) but having an anticanonical divisor $D$ with local canonical and log-canonical thresholds at a point $v$ equal to $(N-2)/N$ and $(N-1)/N$, respectively. This is a claimed special-example obstruction to (1) implying (2), not a statement about sufficiently general hypersurfaces.

The original additionally claims that under (1),(2), $X\times\PP^1$ has only its two projection Mori-fiber structures, and then says every birational map to a Mori fiber space is an isomorphism. The aid flags the latter literal assertion as false: $\operatorname{Bir}_X(X\times\PP^1)=\operatorname{PGL}_2(\C(X))$ in the complex interpretation. Classifying structures up to an appropriate square equivalence differs from prohibiting these fiberwise birational maps. The aid only records the projections as the two regular elementary contractions; it does not establish the full birational-structure classification.
\subsection{Short English statement}
Compare movable-system canonicity with birational superrigidity and all-divisor singularity conditions, especially for quartic threefolds and quintic fourfolds. The product claim also needs an equivalence convention because fiberwise birational maps need not be regular.
\subsection{Sources}
\begin{itemize}
\item Cheltsov, Problem 17 according to supplied verification, AIM \emph{Numerical invariants of singularities and higher-dimensional algebraic varieties}: \url{https://aimath.org/WWN/singularvariety/problems_final.pdf}. Canonical location: \texttt{aim-algebraic-geometry-notes.json}, notes[316]. The extraction instead begins ``7'' and ends with an isolated ``1''; no new source visit is asserted.
\item Original reference: Pukhlikov [8], for sufficiently general degree-$N$ hypersurfaces with $N\ge6$. The [8] bibliography entry and precise canonical/plt version are not supplied.
\item The aid invokes the Noether--Fano criterion and de Fernex's theorem with its erratum without full citations, and proposes the special cone-section examples with low novelty confidence. It is machine generated and not peer reviewed.
\item Local registry: \texttt{data/unsolvedmath/problems.json}, permanent catalogue entry, exact source code \texttt{catalogue entry}.
\end{itemize}
\end{document}
